\documentclass[12pt,a4paper]{article}

\usepackage[T1]{fontenc}
\usepackage{lmodern}
\usepackage[a4paper,left=25mm,right=25mm,top=22mm,bottom=23mm]{geometry}
\usepackage{microtype}
\usepackage[hidelinks]{hyperref}

\hypersetup{
    pdftitle={Statement of Research Interests},
    pdfauthor={Lu Niu},
    pdfsubject={Machine-learning surrogate models for two-dimensional materials}
}
\setlength{\parindent}{0pt}
\setlength{\parskip}{0.55\baselineskip}
\linespread{1.04}
\raggedbottom

\begin{document}

\begin{center}
    {\Large\bfseries Statement of Research Interests}\par
    \vspace{0.2em}
    {\large Lu Niu}\par
    School of Physics, The University of Sydney
\end{center}
\vspace{-0.4em}

My research examines how atomic structure controls the electronic properties of two-dimensional materials.
During my PhD at the University of Sydney, I have investigated this connection through first-principles calculations of graphene-based heterostructures, boron allotropes, and pristine and hydrogen-functionalized beryllium sheets.
I now wish to extend this work to twisted multilayers, where relaxation and disorder demand models that reach beyond the practical scale of direct density functional theory (DFT).
The postdoctoral project in Professor Christoph Ortner's group at UBC offers a direct route towards this aim through machine-learning interatomic potentials (MLIPs) and machine-learning tight-binding (MLTB) models.

My work on graphene--borophene and graphene--boron-carbide bilayers connects stacking registry and interlayer separation to binding, charge redistribution, and electronic and optical response.
I compared competing configurations, optimized their structures, and examined their bands and dielectric functions using PBE and HSE06 calculations in VASP.
This work taught me to resolve small energy differences carefully and to distinguish changes caused by geometry from those introduced by the electronic-structure approximation.
My subsequent studies of boron and beryllium systems extended these comparisons to hydrogen termination, lattice stability, magnetism, and band topology.
Two studies have led to first-author publications in \textit{Nanomaterials} and \textit{Physical Review Materials}; the beryllium manuscript is in preparation.

Alongside these calculations, I have developed reusable analysis scripts and Jupyter workflows for convergence tests, electronic bands, phonons, molecular dynamics, and optical response.
I also have experience with Quantum ESPRESSO, Wannier90, and EPW through electron--phonon calculations performed with collaborators' guidance.
These activities provide a practical foundation for generating DFT reference data and testing whether a model preserves the physical quantities of interest.
My machine-learning experience currently consists of small toy-model experiments in JAX.
I would build on this beginning while contributing my established experience in first-principles modelling and numerical analysis.

At UBC, I would first focus on twisted graphene as a controlled starting point, building on the group's atomic cluster expansion potential for twisted multilayers~\cite{wang2025ace}.
The central question I wish to address is how accurately a structural model must reproduce local stacking and interlayer forces to predict the electronic consequences of moir\'e relaxation.
I would begin by reproducing selected structural benchmarks, then identify configurations that challenge the existing model, such as strained regions, corrugation, and selected defect environments.
This would establish a manageable initial project before extending the approach to BN or transition-metal dichalcogenides in response to the collaboration's priorities.

Next, I would generate and curate additional reference calculations with consistent settings and recorded convergence criteria.
The data would sample changes in stacking, layer separation, and atomic displacement, including configurations away from equilibrium.
I would test the potential against relative registry energies, forces, relaxed geometries, and selected phonon modes, reserving entire families of configurations for validation.
This would expose errors that a random division of closely related structures could conceal.
Comparisons with new DFT calculations would then guide further sampling towards configurations where the model fails or its predictions become uncertain.

The electronic part of the project would follow from these structural tests.
I would develop a separately validated MLTB description, drawing on the group's work on equivariant Hamiltonian models~\cite{zhang2022equivariant}.
A consistent orbital representation, Hermiticity, and the appropriate symmetry transformations would constrain the construction.
For cells accessible to DFT, I would compare band dispersions, orbital character, and the density of states before applying the model to larger relaxed structures.
I am particularly interested in tracing how errors in interlayer spacing or local strain propagate into band widths, small gaps, and electronic localization.
This connects the accuracy of the structural and electronic models to the observables they are intended to predict.

With these comparisons established, I would examine how defects and nonuniform relaxation modify moir\'e domain patterns and the local density of states.
The collaboration's contact with experiment would help select geometries and observables that answer concrete materials questions.
I would aim to deliver reproducible reference datasets, documented models, and clear tests of their range of validity, alongside the physical results.

My longer-term interests include the interplay between electronic structure and many-body phenomena in two dimensions.
Within the exploratory part of the position, I would like to investigate how validated low-energy Hamiltonians could support subsequent studies of interactions in selected moir\'e systems.
The immediate priority would remain the structural and electronic surrogate models that make such questions tractable.
I am therefore particularly interested in combining my computational materials background with the group's mathematical and numerical expertise to develop models whose accuracy can be assessed through their physical predictions.

\vspace{0.1em}
{\small
\begin{thebibliography}{2}
\setlength{\itemsep}{0pt}
\setlength{\parskip}{0pt}
\bibitem{wang2025ace}
Y. Wang et al., ``An Atomic Cluster Expansion Potential for Twisted Multilayer Graphene,''
\href{https://doi.org/10.1088/2632-2153/ae1807}{\textit{Machine Learning: Science and Technology} \textbf{6}, 045040 (2025)}.

\bibitem{zhang2022equivariant}
L. Zhang et al., ``Equivariant analytical mapping of first principles Hamiltonians to accurate and transferable materials models,''
\href{https://doi.org/10.1038/s41524-022-00843-2}{\textit{npj Computational Materials} \textbf{8}, 158 (2022)}.
\end{thebibliography}
}

\end{document}
